\documentclass{article}
\usepackage{/globals/anki}
\ankiprefix{supplychain-systeme}
\ankitopic{Supply-Chain-Managements}
\ankishow
\begin{document}
\section{Supply-Chain-Managements} 
Beim Supply-Chain-Managements wird versucht Probleme innerhalb der Lieferungskette vorrauszusehen und zu umgehen.
\begin{description}
 \item[Cooperative Planning, Forecasting and Replenishment] Eine Prognose, die alle Zwischenlieferanten beinhaltet.
 \item[Available-to-Promise] Das Prüfen von Alternativen.
 \item[Supply-Chain-Event-Management] Störungen vermeiden; benachrichtigen, planen.
 \item[Vendor Managed Inventory] Den Lieferanten Lagerzugriff geben.
\end{description}
\ankinote{cpfr}{Cooperative Planning, Forecasting and Replenishment}{Prognose mit allen Zwischenlieferanten}
\ankinote{atp}{Available-to-Promise}{Alternativen prüfen}
\ankinote{scem}{Supply-Chain-Event-Management}{Störungen vermeiden, benachrichtigen, planen}
\ankinote{vmi}{Vendor Managed Inventory}{Lieferanten Lagerzugriff}
\end{document}